\documentclass[a4paper, 12pt]{article}

\usepackage{utils}

\renewcommand*{\today}{17 March 2026}

\begin{document}

\hotbox{Géometrie 2}{CM 6}{\today}

\section{Les droites}

\subsection{Représentation analytique}

\begin{definition}[(Équation (cartésienne) d'une droite dans un plan)]
    Soit $D$ une droite de $\R^2$. On dit que $D$ admet une \textbf{équation cartésienne} s'il existe $a, b, c \in \R$ où $a$ et $b$ ne sont pas tous les deux nuls, tels que
    $$
        D = \{ (x, y) \in \R^2 : ax + by + c = 0 \}
    $$
\end{definition}

\begin{definition}[(Équation (cartésienne) d'une droite dans l'espace)]
    Une droite $D$ de $\R^3$ est définie par l'intersection de deux plans. Par conséquent, $D$ admet une \textbf{équation cartésienne} s'il existe \text{$n_1 := (u_1, v_1, w_1), n_2 := (u_2, v_2, w_2) \in \R^3$} et $d_1, d_2 \in \R$ où 
    $n_1$ et $n_2$ ne sont pas colinéaires, tels que
    $$
        \begin{cases}
            u_1 x + v_1 y + w_1 z + d_1 = 0 \\
            u_2 x + v_2 y + w_2 z + d_2 = 0
        \end{cases}
    $$
\end{definition}

\begin{definition}[(Équation paramétrique d'une droite)]
    Soit $D$ une droite de $\R^2$. On dit que $D$ admet une \textbf{équation paramétrique} s'il existe un point $A = (x_0, y_0) \in D$ et un vecteur directeur $\vec{u} = (u, v) \in \R^2\setminus\{0\}$ où $u$ et $v$ tels que
    $$
        D = \{ A + \lambda\vec{u} : \lambda \in \R \}
    $$
\end{definition}

\end{document}